\begin{center}
    \section*{\fontsize{20}{20}\selectfont Chapter 2}
\end{center}
\vspace{10mm}
\section{Background Study}

Building a system like EHealth does not happen in a vacuum. There is a growing body of work across three distinct areas that this project sits at the intersection of: using large language models for medical triage and clinical decision support, using blockchain to anchor the integrity of medical records, and using automatic speech recognition (ASR) paired with LLMs to automate clinical documentation. This chapter surveys the most relevant work in each area, identifies what those systems got right, where they fell short, and explains how EHealth's design choices relate to the existing literature.

\subsection{Literature Review}

\subsubsection{LLM-Based Medical Triage and Chatbots}

The question of whether language models can assist with clinical triage has been studied seriously since GPT-3.5 and GPT-4 became broadly accessible. Comparative studies have evaluated LLM and ChatGPT-based triage directly against trained and untrained emergency department staff, finding that model-assisted triage can approach professional performance and can meaningfully raise the performance of untrained personnel \cite{triage_llm_2024}. This is a significant result because it reframes the goal: the model does not need to replace a clinician, it just needs to be better than no guidance at all.

\noindent Hallucination is the most cited risk in applying LLMs to medical tasks. One response to this is retrieval-augmented generation (RAG), which grounds the model's output in a curated knowledge base rather than relying on parametric recall. Retrieval-augmented designs have been proposed specifically to reduce hallucination in preliminary diagnosis chatbots \cite{inclusive_healthcare_2025}. EHealth does not use RAG, primarily because the triage task is constrained: the model is not being asked to reason over clinical literature, only to map a symptom description to one of three urgency levels. The constrained JSON output schema acts as a guardrail that serves a similar stabilizing function.

\noindent Layered architectures have also been described for general healthcare chatbots, separating the LLM reasoning core, a response formatting module, and a safety or filtering layer \cite{healthcare_chatbot_2025}. EHealth follows an analogous pattern: Baichuan-M2-32B handles reasoning, the prompt template enforces JSON schema, and a service layer handles non-medical input as a separate branch that never reaches the triage classification logic. Broader surveys catalogue the wider risks of deploying LLMs as medical chatbots at scale --- hallucination, demographic bias, the need for human oversight --- and highlight that no current system should be positioned as a diagnostic replacement \cite{llm_chatbot_risks_2025}. EHealth is explicit on this point: the triage output is a first-pass urgency estimate, not a diagnosis, and the platform's design reflects that framing.

\subsubsection{Blockchain-Anchored Medical Record Integrity}

A recurring pattern across the blockchain-for-healthcare literature is to keep bulk data off-chain and anchor only a cryptographic hash on-chain. This keeps the ledger tamper-evident without the cost or privacy exposure of storing clinical content directly \cite{medical_blockchain_ehr_2022, blockchain_healthcare_2020}. The rationale is straightforward: blockchain storage is expensive compared to conventional databases, and putting medical records on a public chain creates a privacy problem that is difficult to solve. Anchoring only a hash avoids both issues while preserving the core property: once a hash is on-chain, any modification to the underlying file produces a different hash, and the mismatch is immediately detectable.

\noindent Ethereum and IPFS-based designs have explored combining off-chain content-addressed storage with smart-contract-mediated access control for patient-centric record sharing \cite{blockchain_record_sharing_2025}. More recent frameworks extend this to soulbound tokens and access-control policies enforced entirely in the contract \cite{blockchain_soulbound_2024}. EHealth adopts the simpler hash-anchoring pattern for a narrower purpose: verifying a single artifact (a prescription PDF) rather than managing a full longitudinal record. This keeps the on-chain footprint minimal --- a fixed 32-byte keccak256 digest per prescription --- and keeps gas cost predictable regardless of prescription length or complexity. The dual-check response (database match and on-chain match independently reported) gives the patient an honest, server-agnostic verification path rather than a single boolean they have to take on trust.

\subsubsection{ASR and LLMs for Clinical Documentation}

Whisper-family models are now the dominant open ASR backbone in clinical documentation pipelines. The standard pattern is a two-stage pipeline: Whisper handles transcription, and an LLM restructures the raw transcript into a clinical note or form \cite{asr_healthcare_2026}. Proof-of-concept systems have shown that a Whisper-plus-LLM pipeline can populate structured clinical forms from doctor--patient audio in well under a minute for several-minute consultations \cite{llm_transcription_2024}. Other work benchmarks ASR-coupled LLM modules specifically for medical diagnostics under noisy, device-heterogeneous recording conditions \cite{asr_llm_benchmarking_2025}.

\noindent EHealth follows the same two-stage pattern but makes one significant architectural choice that distinguishes it from most published work: the entire pipeline runs inside the deployment's own containers. No audio or transcript is sent to a third-party ASR API. Xenova/whisper-tiny runs via the Hugging Face Transformers.js library on the ONNX runtime, and the LLM call goes to Baichuan-M2-32B only after local transcription is complete. This matters for two reasons. First, it means patient audio never leaves the server environment. Second, it means the pipeline's latency and cost are predictable and not subject to third-party rate limits. The trade-off is accuracy: whisper-tiny is the smallest Whisper checkpoint, and a larger model would produce lower word error rates, especially on medical terminology. At 14.2\% general WER and 8.6\% medical WER on 50 test dialogues, the accuracy is practically useful, but larger checkpoints would likely improve particularly on uncommon drug names and specialist vocabulary.

\subsubsection{Summary Comparison Table}

Table~\ref{tab:lit_review} summarizes the key related works, their domain focus, approach, and how EHealth relates to or differs from each.

\vspace{3mm}
{
\renewcommand{\arraystretch}{1.3}
\begin{xltabular}{\textwidth}{|>{\centering\arraybackslash}p{0.8cm}|>{\raggedright\arraybackslash}p{2.0cm}|>{\raggedright\arraybackslash}p{2.2cm}|>{\raggedright\arraybackslash}X|>{\raggedright\arraybackslash}p{3.5cm}|}

% 1. FIRST HEAD (First page header)
\hline
\textbf{Ref} & \textbf{Focus} & \textbf{Approach} & \textbf{Key Finding / Gap} & \textbf{Relation to EHealth} \\
\hline
\endfirsthead

% 2. HEAD (Continuation pages header)
\hline
\textbf{Ref} & \textbf{Focus} & \textbf{Approach} & \textbf{Key Finding / Gap} & \textbf{Relation to EHealth} \\
\hline
\endhead


% 4. LAST FOOT (only appears once, at the very end of the table)
\hline
\caption{Summary Overview of Literature Review} \label{tab:lit_review} \\
\endlastfoot


% 5. TABLE DATA
\cite{triage_llm_2024} & LLM triage vs. clinical staff & ChatGPT vs. ED triage nurses & LLM-assisted triage approaches professional accuracy; raises untrained performance & EHealth's constrained prompting adopts same framing; no diagnostic claim \\
\hline
\cite{inclusive_healthcare_2025} & Preliminary diagnosis chatbot & Retrieval-augmented LLM & RAG reduces hallucination vs. direct LLM query & EHealth uses schema-constrained output instead of RAG for triage \\
\hline
\cite{healthcare_chatbot_2025} & General healthcare chatbot & Layered LLM architecture & Separating reasoning, formatting, and safety layers improves reliability & EHealth mirrors this: LLM + JSON schema + non-medical branch separation \\
\hline
\cite{llm_chatbot_risks_2025} & LLMs in medical chatbots (survey) & Literature survey & Hallucination, bias, and lack of oversight are systemic risks & EHealth positions triage as urgency estimate, not diagnosis \\
\hline
\cite{medical_blockchain_ehr_2022} & EHR integrity on blockchain & Off-chain data + on-chain hash & Hash-anchoring gives tamper-evidence without privacy exposure & EHealth uses same pattern: SHA-256 hash $\rightarrow$ keccak256 $\rightarrow$ Polygon PoS \\
\hline
\cite{blockchain_healthcare_2020} & Smart contracts for healthcare & Ethereum smart contracts & Decentralized record management; high gas cost on Ethereum & EHealth uses Polygon PoS to reduce cost to $\approx \$0.0005$/write \\
\hline
\cite{blockchain_record_sharing_2025} & Patient-centric record sharing & Ethereum + IPFS & Content-addressed off-chain storage + on-chain access control & EHealth stores PDFs on Google Drive; verification only is on-chain \\
\hline
\cite{blockchain_soulbound_2024} & Medical record verification & Soulbound tokens + cloud & Access control fully in smart contract; limited to permissioned actors & EHealth uses open verify endpoint; patients self-verify with file upload \\
\hline
\cite{asr_healthcare_2026} & ASR in post-LLM healthcare (review) & Scoping review & Whisper-family models dominate; larger checkpoints vs. latency tradeoff & EHealth uses whisper-tiny for local, no-network inference \\
\hline
\cite{llm_transcription_2024} & Doctor--patient audio summarization & Whisper + GPT summarization & Structured forms from audio in under one minute on several-minute calls & EHealth achieves same; runs entirely in-container, no API dependency \\
\hline
\cite{asr_llm_benchmarking_2025} & ASR + LLM for medical diagnostics & ASR-LLM pipelines & Device heterogeneity and noise substantially affect WER in real settings & EHealth normalizes to 16\,kHz mono WAV via FFmpeg before transcription \\
\hline

\end{xltabular}
}
% \vspace{2mm}
\subsection{Problem Analysis}

The literature survey makes it clear that none of the reviewed works combine all three problem areas into a single integrated patient-facing product. Most published systems are either purely demonstrative (a chatbot that triages but has no booking or consultation capability), purely infrastructure-focused (a blockchain scheme for records that assumes clinical data is generated elsewhere), or purely pipeline demonstrations (Whisper-plus-LLM documentation without any delivery mechanism for the output).

\noindent This fragmentation is not accidental --- it reflects how academic research is structured. Each paper solves one problem cleanly and evaluates it in isolation. But for a real patient, the value is in the integration. A triage result that does not connect to a booking flow is just information with no action path. A consultation with no automatic documentation means the doctor still does all the paperwork manually. A prescription with no tamper-evidence is just a PDF that anyone could edit.

\noindent Three specific gaps in the existing work are worth naming explicitly. First, there is no published system that connects LLM-based urgency triage directly to a live video consultation booking within the same application session. Second, local (in-container) ASR-based documentation has been demonstrated, but not as part of a system that also handles payment, WebRTC conferencing, and downstream prescription management. Third, blockchain-based prescription integrity has been explored conceptually, but detailed cost measurement on a live public network (as distinct from a private test chain) is rarely reported.

\noindent EHealth addresses all three of these gaps --- not by being more sophisticated than any individual published system in its domain, but by being the first to integrate all three into a deployable, evaluated end-to-end platform.

\vspace{5mm}

\noindent\textbf{Summary:} This chapter surveyed eleven works across three research areas relevant to EHealth: LLM-based medical triage and clinical chatbots, blockchain-anchored medical record integrity, and ASR-driven clinical documentation. The summary comparison table (Table~\ref{tab:lit_review}) shows that while each component area is reasonably mature, no existing work addresses all three simultaneously within a single integrated patient journey. The problem analysis identifies three specific integration gaps --- triage-to-booking continuity, in-container ASR-documentation within a full telemedicine workflow, and live on-chain cost measurement for prescription integrity --- that EHealth fills. The following chapters describe how the system was designed and built to address these gaps.